\documentclass[11pt, letterpaper]{article}

\usepackage[margin=0.75in, top=0.7in, bottom=0.7in]{geometry}
\usepackage{amsmath, amssymb}
\usepackage{booktabs}
\usepackage{graphicx}
\usepackage{hyperref}
\usepackage{float}
\usepackage{caption}
\usepackage{enumitem}
\usepackage{parskip}
\usepackage[T1]{fontenc}
\usepackage{lmodern}
\usepackage{microtype}
\usepackage{fancyhdr}
\usepackage{xcolor}
\usepackage{cmbright}
\usepackage{multicol}
\hypersetup{
    colorlinks=true,
    linkcolor=blue!60!black,
    citecolor=blue!60!black,
    urlcolor=blue!60!black,
}

\pagestyle{fancy}
\fancyhf{}
\rhead{\small Quantathon 2026}
\lhead{\small Columbus LEAP Funding Model --- Executive Summary}
\cfoot{\thepage}

\title{
    \vspace{-1.5cm}
    \textbf{Columbus LEAP Funding Model} \\[0.2cm]
    \large Executive Summary \(\mid\) Quantathon 2026 --- SIAM-MTI Challenge \\[0.1cm]
    \normalsize Quantathon 2026 --- SIAM-MTI Challenge
}
\author{}
\date{}

\begin{document}

\maketitle
\thispagestyle{fancy}
\vspace{-0.5cm}

%--------------------------------------------------------------
\section*{The Question}

How much should the City of Columbus reserve in 2026 so that its Lead Exposure Assistance Program (LEAP)---a zero-interest, deferred-repayment loan fund for water service line replacement---stays solvent through the 2037 federal deadline?

\vspace{-0.2cm}
\section*{The Answer}

\begin{center}
\renewcommand{\arraystretch}{1.3}
\begin{tabular}{@{}lccc@{}}
\toprule
\textbf{Model} & \textbf{95\%-Safe Reserve} & \textbf{Expected} & \textbf{Grant Equivalent} \\
\midrule
Flat-Rate Monte Carlo (3{,}000 paths) & \textbf{\$2.73M} & \$2.49M & \$4.01M \\
ML-Enhanced Agent Simulation (1{,}000 paths) & \textbf{\$3.60M} & \$3.17M & \$6.34M \\
\bottomrule
\end{tabular}
\end{center}

\noindent \textbf{Primary recommendation:} Reserve \textbf{\$2.73M--\$3.60M} under the loan structure, depending on risk tolerance. The loan design saves \textbf{\$1.27M--\$2.74M} versus a grant alternative by recycling repayments.

%--------------------------------------------------------------
\vspace{-0.2cm}
\section*{How We Got There: 8 Methods Converge}

We track 34{,}269 eligible lines across 12 years via cohort-based Monte Carlo simulation. Eight independent methods converge: \textbf{(1)} MC simulation (3{,}000 paths, P95 = \$2.73M); \textbf{(2)} deterministic Markov benchmark (\$2.47M, within 0.7\%); \textbf{(3)} bootstrap inference (P95 CI width: \$24{,}828); \textbf{(4)} tornado sensitivity (uptake rate = 78\% of variance); \textbf{(5)} NPV at 3\%: \$2.05M, stable $\pm$3\% across 1--5\%; \textbf{(6)} 3 neural networks ($r{=}0.997$, $r{=}0.973$) for per-parcel heterogeneity; \textbf{(7)} ML-enhanced agent simulation with 7 enhancements (regime switching, Weibull hazard, urgency ramp, appreciation, contractor cap, fund returns, LGD); \textbf{(8)} Columbus v4 out-of-sample testing. \textbf{CV: 0.11\%} --- extremely stable.

%--------------------------------------------------------------
\vspace{-0.2cm}
\section*{Key Assumptions (All Data-Grounded)}

\begin{center}
\renewcommand{\arraystretch}{1.15}
\small
\begin{tabular}{@{}llll@{}}
\toprule
\textbf{Parameter} & \textbf{Base} & \textbf{Range} & \textbf{Source} \\
\midrule
Eligible lines & 34{,}269 & --- & System-materials table (lead + galvanized, net overlap) \\
Uptake rate & 0.20\%/yr & 0.16--0.25\% & 62 actual loans / 34{,}269 pool = 0.18\%; rounded up \\
Cost inflation & 3\%/yr & 2--4\% & CPI-Urban + construction labor index; \$10K cap \\
Sale probability & 5.6\%/yr & 4.5--6.5\% & Columbus tenure data: avg 18 yrs $\Rightarrow$ 1/18 = 5.6\% \\
Equity extraction & 2.0\%/yr & 1.5--2.8\% & CFPB cash-out refinance data \\
Fund return & 4\%/yr & 3.5--4.5\% & Municipal bond yield benchmark \\
Contractor cap & 120--250/yr & 100--300 & LEAP program manager estimates \\
\bottomrule
\end{tabular}
\end{center}

%--------------------------------------------------------------
\vspace{-0.2cm}
\section*{Stress Testing: The Model Survives Extremes}

\begin{center}
\renewcommand{\arraystretch}{1.2}
\begin{tabular}{@{}lccc@{}}
\toprule
\textbf{Scenario} & \textbf{P95} & \textbf{vs.\ Base} & \textbf{Key Shocks} \\
\midrule
Base Case & \$2.73M & --- & Standard assumptions \\
2029 Housing Crash & \$3.14M & +14.9\% & Sales $-$40\%, equity $-$50\% \\
Surge Uptake & \$4.98M & +82.1\% & 2$\times$ trigger rate, +1\% inflation \\
Perfect Storm & \$5.75M & +110.7\% & Both combined \\
\bottomrule
\end{tabular}
\end{center}

%--------------------------------------------------------------
\vspace{-0.2cm}
\section*{Year-by-Year Cash-Flow Profile}

Outflows peak early (\$719K in 2026) as the at-risk pool is largest. Repayments build gradually; annual inflows exceed outflows starting in \textbf{2035}. By 2037, 42\% of capital is recovered; the Markov chain projects \textbf{80\% recovery by 2051} and \textbf{90\% by 2059}.

%--------------------------------------------------------------
\vspace{-0.2cm}
\section*{Geographic Risk Concentration}

Composite risk scoring (35\% lead risk + 25\% inverse property value + 20\% eligible share + 10\% renter rate + 10\% cost) reveals that zip codes \textbf{43201} (score 0.92, 98\% eligible) and \textbf{43206} (score 0.84, 99\% eligible) carry the highest demand risk. These areas should be prioritized for outreach and monitoring.

%--------------------------------------------------------------
\vspace{-0.2cm}
\section*{Decision Support: Adaptive Strategy \& Monitoring}

Rather than committing the full reserve upfront, Columbus can adopt a \textbf{two-stage adaptive strategy}:

\begin{itemize}[nosep, leftmargin=1.5em]
    \item \textbf{Stage 1:} Reserve \$1.64M in 2026 (covers expected demand through 2029).
    \item \textbf{Stage 2:} If cumulative originations exceed the P75 threshold by 2029, top up by \$1.09M.
    \item \textbf{Expected total:} \$2.73M --- same as the flat recommendation, but with optionality.
\end{itemize}

\noindent \textbf{Three early-warning triggers} signal when to reassess:

\begin{center}
\renewcommand{\arraystretch}{1.2}
\small
\begin{tabular}{@{}lll@{}}
\toprule
\textbf{KPI} & \textbf{Warning Threshold} & \textbf{Action} \\
\midrule
Cumulative originations by 2028 & $>$ 193 loans & Review reserve adequacy \\
Average funded amount & $>$ \$8{,}312 & Increase inflation assumption \\
Cumulative repayments by 2029 & $<$ 23 & Review repayment assumptions \\
\bottomrule
\end{tabular}
\end{center}

%--------------------------------------------------------------
\vspace{-0.2cm}
\section*{What Would Reduce Uncertainty Most}

Value-of-information analysis shows that \textbf{78\% of model variance} comes from LEAP uptake rate uncertainty. A multi-year uptake history would shrink the confidence interval by \$19{,}391. The break-even uptake rate is \textbf{0.202\%/year}---below this, the fund stays solvent at any reasonable funding level.

%--------------------------------------------------------------
\vspace{-0.2cm}
\section*{Financial Reasonableness Check}

\begin{multicols}{4}
\centering
\textbf{$\sim$\$80/line} \\ Reserve per eligible line

\columnbreak
\textbf{\$7{,}386 avg} \\ Observed per-loan cost

\columnbreak
\textbf{42\% by 2037} \\ Principal recovered

\columnbreak
\textbf{$<$0.3\% budget} \\ Of Columbus annual budget
\end{multicols}

\vspace{-0.2cm}
\noindent\rule{\textwidth}{0.4pt}
\vspace{0.1cm}

\noindent \textbf{Bottom line:} Eight independent methods, three stress scenarios, and five discount rates all confirm that a \textbf{\$2.73M loan-based reserve} (95\%-safe, base case) is sufficient to operate LEAP through 2037. The loan structure saves \$1.27M over grants. The recommendation is robust, validated, and actionable.

\end{document}
